%% ==================== 参考文献 模板 (v1.5.6, BZD 数模社 2026 规范) ====================
%%
%% 借鉴 BZD 数模社 参考文献模板: 5-15 条, GB/T 7714 格式
%%
%% BZD 原则:
%% - 顺序编码制, 按正文第一次引用的顺序编号为 [1] [2] [3]...
%% - 真实可核验, AI 推荐的文献必须由人工核对作者/题名/期刊/年份
%% - AI 工具 (ChatGPT/DeepSeek/Copilot) 禁止列入参考文献
%% - 5-15 条最佳, 优先经典 + 近 3-5 年 + 数据来源 + 标准
%% - 每条格式严格: 作者. 题名 [J/M/D/C/N/R/S/P/EB/OL/DB/OL/Z]. 出版地: 出版者, 年.
%%
%% 跑题用户必须:
%% 1. 删掉所有 % 注释
%% 2. 替换【】占位符为自己的实际参考文献
%% 3. 删 下面的 \RefSlot{...} 行 (check 15 抓这个 inline 【】)
%% 4. 按正文实际引用顺序重新编号

\section{参考文献}

% ====== 模板示例 (按 BZD 通用模板, 跑题后必须替换为真实文献) ======

\begin{enumerate}[label={[\arabic*]}, leftmargin=2em]
  \item \textbf{【作者 1】.} \textbf{【论文题名】} [J]. \textbf{【期刊名】}, \textbf{【年份】}, \textbf{【卷(期)】}: \textbf{【起止页码】}.
  \item \textbf{【作者 2】.} \textbf{【书名】} [M]. \textbf{【出版地】}: \textbf{【出版社】}, \textbf{【出版年】}: \textbf{【起止页码】}.
  \item \textbf{【作者 3】.} \textbf{【论文题名】} [D]. \textbf{【出版地】}: \textbf{【授予单位】}, \textbf{【出版年】}.
  \item \textbf{【发布机构】.} \textbf{【报告名称】} [R]. \textbf{【出版地】}: \textbf{【发布机构】}, \textbf{【出版年】}.
  \item \textbf{【发布机构】.} \textbf{【数据集名称】} [DB/OL]. \textbf{【网址】}, \textbf{【更新时间】} / \textbf{【访问日期】}.
  \item \textbf{【作者或机构】.} \textbf{【网页资源名称】} [EB/OL]. \textbf{【网址】}, \textbf{【发布日期】} / \textbf{【访问日期】}.
  \item \textbf{【标准号】.} \textbf{【标准名称】} [S]. \textbf{【出版地】}: \textbf{【出版者】}, \textbf{【出版年】}.
  \item \textbf{【颁布机构】.} \textbf{【文件名称】} [Z]. \textbf{【发布机构】}, \textbf{【发布日期】}.
\end{enumerate}

% 禁止: ChatGPT / DeepSeek / Claude / Copilot / 文心一言 / 通义千问 等 AI 工具条目.
% 禁止: AI 虚构作者 / 题名 / DOI / 期刊 / 页码 / 网址.
% 高风险: 正文有引用编号, 但文末没有对应条目 (一对一对应).
% 推荐: 经典模型原始文献 (Breiman 2001 Random Forest, Lo 1997 NIPT, Chiu 2008...) + 近 5 年相关 + 官方数据/标准.

% 跑题后必须删下面这一行 (check 15 占位符自检)
\RefSlot{【这里填 5-15 条真实可核验的参考文献, 按正文引用顺序编号, 删掉这一行】}
